\label{ch:methodology}

\subsection{Evaluation goals}

The experiments answer three operational questions:

\begin{enumerate}
  \item Does enabling migration batching improve user-visible serving metrics
  under asynchronous EPLB?
  \item Does it reduce high-percentile network traffic, the mechanism the
  scheduler is designed to affect?
  \item Is the CPU scheduling overhead small relative to the observed E2E time
  saved?
\end{enumerate}

These questions lead to three directional expectations rather than formal
statistical hypotheses. First, batching should lower the high tail of observed
NIC traffic because fewer migration flows contend at one rank. Second, if that
reduction protects foreground communication, output throughput should increase
and TTFT, TPOT, and E2E latency should decrease. Third, cycle-wide schedule
construction should consume only a small fraction of the serving time saved.
The single A/B pair can check whether the observations have these directions,
but it cannot establish significance or estimate run-to-run variance.

Two separate runs are used for two separate purposes. A clean A/B serving run
compares batching off and on without profiler instrumentation. A profiling run
repeats the batching-on condition with NVTX-only Nsight Systems capture. This
separation prevents profiler overhead from contaminating the reported serving
comparison.

\subsection{Hardware and software environment}

\begin{table}[t]
  \centering
  \caption{Experimental environment.}
  \label{tab:environment}
  \begin{tabularx}{\textwidth}{lX}
    \toprule
    Component & Configuration \\
    \midrule
    Nodes & Four Ubuntu 22.04.5 servers \\
    Accelerators & One NVIDIA Quadro RTX 6000, \SI{24}{GB}, per node \\
    Network & \SI{10}{Gb/s} Ethernet, interface \code{eno2}, no RDMA \\
    Model & Qwen/Qwen3-30B-A3B-Instruct-2507, FP16 \\
    Parallelism & Tensor parallel size 4; expert parallel enabled \\
    Runtime & Experimental vLLM build; Ray 2.56.1 \\
    Migration backend & NIXL 1.3.2 over UCX/TCP \\
    Execution & Eager mode; prefix caching enabled \\
    Profiling & NVIDIA Nsight Systems 2024.5.1, NVTX-only \\
    \bottomrule
  \end{tabularx}
\end{table}

The four GPU processes serve as both TP and EP ranks. One GPU per host makes
each inter-rank expert transfer an inter-node transfer, avoiding ambiguity
between local NVLink/PCIe movement and Ethernet traffic. The hardware is not a
modern production cluster; its relatively constrained network is useful for
observing contention but limits external validity.

NCCL and Gloo are bound to \code{eno2}, InfiniBand is disabled, and UCX is
also bound to the same device. NIXL is selected explicitly rather than relying
on automatic backend selection. Both A/B conditions use the same software
revision and network configuration.

\subsection{Model and EPLB configuration}

The evaluated model contains 48 layers and 128 experts, selecting eight experts
per token \cite{qwen3modelcard}. Server maximum sequence length is 4096 and GPU
memory utilization is set to 0.9. EPLB uses the configuration in
\cref{tab:eplb-config}.

\begin{table}[t]
  \centering
  \caption{EPLB parameters common to the A/B experiment.}
  \label{tab:eplb-config}
  \begin{tabular}{lr}
    \toprule
    Parameter & Value \\
    \midrule
    \code{window\_size} & 50 \\
    \code{step\_interval} & 50 \\
    \code{num\_redundant\_experts} & 32 \\
    \code{log\_balancedness} & false \\
    \code{use\_async} & true \\
    \code{communicator} & \code{nixl} \\
    \code{enable\_migration\_batching} & false / true \\
    \bottomrule
  \end{tabular}
\end{table}

The short observation and step intervals deliberately cause recurring
rebalances during the finite benchmark. This is a stress-oriented configuration
for studying the migration path rather than a recommended production default.
The 32 redundant experts create physical flexibility for EPLB to react to
recent routing demand.

\subsection{Workload}

Each condition starts a fresh server and executes an eight-request warm-up.
The measured workload uses \code{vllm bench serve} with 500 successful requests
and maximum concurrency 32. Every request has a 300-token random prefix, 200
additional input tokens, and 300 requested output tokens. The random seed is
zero, temperature is zero, and end-of-sequence tokens are ignored so all
requests produce the configured output length. Prefix caching remains enabled,
matching the intended online configuration rather than disabling cache to
manufacture compute load.

The request set is synthetic. It controls sequence lengths and makes the A/B
conditions comparable, but random tokens are not representative of semantic
production traffic. The experiment studies system interference under frequent
migration, not application-level output quality. Because placement batching
does not modify model computation or routing decisions, no accuracy evaluation
is required for the feature itself.

\subsection{A/B procedure and controls}

The independent variable is
\code{enable\_migration\_batching}. All other documented parameters are held
constant. The off case ran first and the on case second, each with a fresh
server. Both completed 500 of 500 requests with zero failures. The benchmark
client records aggregate duration and request-level latency distributions.

The experiment is a single repetition. It therefore cannot quantify run-to-run
variance, confidence intervals, or statistical significance. Percent changes
are descriptive comparisons of these archived runs. A stronger study would
randomize condition order and repeat each condition enough times to estimate
variance. This limitation is stated explicitly rather than obscured by the
large number of requests, because requests within one long run are not
independent replications of server startup and cluster state.

\subsection{Serving metrics}

For a metric where larger is better, such as throughput, improvement is

\begin{equation}
  I_{\uparrow}=\frac{V_{\mathrm{on}}-V_{\mathrm{off}}}
  {V_{\mathrm{off}}}\times100\%.
\end{equation}

For a latency metric where smaller is better,

\begin{equation}
  I_{\downarrow}=\frac{V_{\mathrm{off}}-V_{\mathrm{on}}}
  {V_{\mathrm{off}}}\times100\%.
\end{equation}

The thesis reports output tokens per second and P50/P99 TTFT, TPOT, and E2E
latency. Request throughput is retained in the raw JSON and the supplementary
table. Because all 500 requests produce exactly 300 output tokens, request and
output throughput necessarily have the same relative change; the main text
uses output throughput to avoid presenting this one result twice. P99 values
are computed by the benchmark client from the 500 request samples.

\subsection{NIC measurement}

NIC traffic is not produced by \code{vllm bench serve}. A separate sampler on
the head node reads the cumulative receive and transmit byte counters under
\path{/sys/class/net/eno2/statistics/} once per second. For adjacent samples at
times $t_{i-1}$ and $t_i$, aggregate traffic is

\begin{equation}
  B_i=\frac{(RX_i-RX_{i-1})+(TX_i-TX_{i-1})}
  {t_i-t_{i-1}}.
\end{equation}

Values are converted to MB/s and the P99 is computed over the formal benchmark
window. Summing RX and TX describes total head-node interface pressure; it is
not physical link utilization and it cannot attribute bytes to EPLB versus
inference. The A/B comparison and time-series shape provide supporting evidence,
not packet-level causal attribution.

\subsection{Scheduler profiling}

The profiling run uses the same model, EPLB settings, workload lengths,
concurrency, and 500 prompts as the batching-on E2E condition. The only change
is that Ray processes are launched in NVTX-only Nsight Systems sessions with
CPU sampling and context-switch tracing disabled. NVTX ranges are designed to
mark application phases with low capture volume \cite{nsightguide}.

The range \code{eplb: schedule migration batches} encloses the call that builds
schedules for all 48 layers. Each rank writes its own \code{rank*.nsys-rep}.
Explicit formal-start and formal-end timestamps exclude model loading and
warm-up. The reports were analyzed with Nsight Systems 2024.5.1. For each rank,
the number of calls, P50 and P99 elapsed duration per call, and cumulative
duration are extracted.

The explicit CUDA NVTX range was temporary measurement instrumentation added
only to make the asynchronous EPLB worker visible to Nsight Systems. It was
used in the profiling run, was absent from the clean E2E A/B run, and was not
included in the final production patch. The scheduler call and placement
inputs were otherwise unchanged.

Ranks construct their identical schedules concurrently on different host
CPUs. Their costs must therefore not be summed to estimate critical-path
overhead. A conservative comparison uses the largest cumulative rank cost,
$C_{\max}$, and clean A/B time saved,

\begin{equation}
  S=T_{\mathrm{off}}-T_{\mathrm{on}}.
\end{equation}

The cost-to-saved ratio is

\begin{equation}
  R=\frac{C_{\max}/1000}{S}\times100\%.
  \label{eq:cost-ratio}
\end{equation}

This ratio does not claim that all scheduling time lies on the serving critical
path. It intentionally compares a directly measured conservative CPU elapsed
cost with the directly measured wall-clock difference.

\subsection{Data provenance and reproducibility}

The authoritative E2E source directory is
\path{results/serving_nixl_20260911_async_random500}. It contains the
unmodified off/on \code{bench.json}, server and client logs, EPLB
configuration, one-second \code{nic.tsv}, warm-up output, plotting source, and
summary CSV files. The authoritative profiling directory is the same name with
suffix \code{\_profile}; it contains four Nsight reports, formal-window
timestamps, benchmark output, exact commands, and \code{profile\_summary.csv}.

Older directories represent intermediate designs or exploratory workloads and
are excluded from all final tables. Reproduction commands are transcribed in
\cref{app:commands}. This explicit provenance prevents accidental selection of
a favorable intermediate run.
